\section{问题四：固定日程下的独立分组与资源配置}
\label{sec:q4}

\subsection{问题分析与求解思路}
\subsubsection{固定任务日程与组间不共享约束}
本问不再优化运输或中继日程，而是研究同一救援预案划分为独立执行单元后的配置代价。固定问题三工期优先方案的箱组、机型、路线、绝对时刻、悬停位置、服务窗口和实际保障关系，改变服务区归属与组内同型资源指派；组间不得共享实体。原来可以跨区域接续使用的一架飞机或一组电池，独立执行时可能需要分别配置。

同型资源允许重新指派，因此不能把问题三已有的实体编号跨组使用直接判为分组失败；需要计算每组在固定日程下最少需要多少同型资源。反之，也不能为减少需求而平移任务时刻，这会改变本问的冻结条件。

\subsubsection{运输与通信关联形成的不可拆分关系}
以服务区为节点，同一运输架次访问的区域必须同组。同一逻辑中继任务实际保障的运输架次所关联区域也必须同组，关联采用正式区间与临界端点的实际提供方，不用理论覆盖范围推测。对此关系求连通闭包，得到三个不可拆块：
\begin{equation}
 B_1=\{S001\},\quad B_3=\{S006\},\quad
 B_2=\mathcal S\setminus(B_1\cup B_3).
\label{eq:q4_blocks}
\end{equation}
$B_2$包含其余13区，不是按地理距离聚成的一类。两组非空无标签分区只有3种，三组只有1种，可完整枚举而不必引入随机聚类算法。

\subsection{分组与分类型资源配置模型}
\subsubsection{合法分组与完整任务归属}
令$u_{bg}\in\{0,1\}$表示任务块$b$归组$g$，满足$\sum_gu_{bg}=1$和各组非空。运输与中继任务按关联块完整归属，不能跨组拆箱、截取中继服务时间或复制同一逻辑任务。资源类型按固定顺序记为
\begin{equation}
 \mathcal P=\{U_A,U_B,U_C,B_A,B_B,B_C,U_R,K_R\},\qquad
 \boldsymbol I=(4,2,2,6,4,4,2,6).
\label{eq:q4_inventory}
\end{equation}
分别对应三类运输机、三类电池、中继机和中继能源组件。类型间不可替代，总件数不代表同等能力或成本。

\subsubsection{固定占用区间与组内最少资源数}
任务$j$对类型$p$的占用为$I_{jp}=[s_{jp},r_{jp})$，释放端点沿用原计划：运输机计至返航，运输电池计至充满，中继机计至返航周转结束，中继组件计至恢复。组$g$的并发需求及最少配置为
\begin{equation}
 n_{gp}(t)=\sum_{j\in\mathcal J_g}\mathbf1\{t\in I_{jp}\},\qquad
 N_{gp}^{\min}=\max_tn_{gp}(t).
\label{eq:q4_peak}
\end{equation}
峰值是必要下界；在当前同型可替换、无额外转移或配对约束的条件下，按开始时刻排序，将每项任务指派给已释放的资源，若没有则新开一个资源，可用恰好峰值数量完成。若必须新开资源，则该时刻所有已有资源都在占用，说明峰值同步增加，故构造达到下界。该性质与区间图的着色结构对应\cite{fulkerson1965interval}，不适用于另有飞行员、共享充电器或跨基地转移限制的扩展模型。

\subsubsection{库存缺口、配置规模与工作量均衡}
各组独立需求与正缺口为
\begin{equation}
 D_p=\sum_gN_{gp}^{\min},\qquad G_p=\max(D_p-I_p,0),\qquad
 G_\Sigma=\sum_pG_p,
\label{eq:q4_gap}
\end{equation}
总配置$D_\Sigma=\sum_pD_p$仅是不同资源件数的描述性总和。库存剩余$R_p=\max(I_p-D_p,0)$不能跨类型抵消缺口，也不是时间闲置率。

组工作量$W_g$按其运输和中继任务从准备开始至返航的时长相加，不包括电池充电与中继返航周转。令$\overline W=K^{-1}\sum_gW_g$，工作量变异系数为
\begin{equation}
 CV_W=\frac{\sqrt{K^{-1}\sum_g(W_g-\overline W)^2}}{\overline W}.
\label{eq:q4_cv}
\end{equation}
按$(G_\Sigma,D_\Sigma,CV_W)$依次比较分区，先补足执行能力，再比较配置与均衡。不是将资源量和工作量任意加权为一个指标。

\subsection{合法分区的完整枚举}
\subsubsection{不可拆任务块压缩与分区生成}
用规范化块标签生成无标签分区，避免把组号交换当成不同方案。对每个候选，将任务映射到唯一所属组，逐资源类型遍历所有开始和释放事件，计算峰值。保留每个峰值的真实时刻、同时占用任务和构造出的资源链，便于核对结果而非只输出总数。

三种两组分区的完整结果见表\ref{tab:q4_partitions}。本轮从正式运输、通信区间及端点记录重新构造，得到与锁定摘要一致的三个块及各类需求，而不是依据已知最优分区反推其他两种结果。
\begin{table}[htbp]\centering\small
\caption{全部三种两组分区的配置与均衡结果}\label{tab:q4_partitions}
\begin{tabular}{@{}llrrr@{}}
\toprule
分区 & 需求向量 & 总配置 & 总缺口 & 工作量CV\\
\midrule
$B_1+B_2\mid B_3$ & 4,2,3,6,4,4,2,3 & 28 & 1 & 0.941762\\
$B_1+B_3\mid B_2$ & 5,2,4,7,4,6,2,3 & 33 & 6 & 0.781197\\
$B_1\mid B_2+B_3$ & 5,2,4,7,4,6,2,3 & 33 & 6 & 0.839435\\
\bottomrule
\end{tabular}
\end{table}


\subsubsection{资源峰值、实体指派与最少配置核验}
所有事件采用未舍入的十进制时间，同一时刻先释放再占用，保持半开区间语义。对每个组和类型，除扫描与贪心构造外，再构建任务可接续有向无环图：若前任务释放不晚于后任务开始则连边。其二部图最大匹配给出最小路径覆盖数$|\mathcal J_{gp}|-|M_{gp}^{\max}|$，与峰值及资源链数量分别核对。

本轮对工期优先、节能、增强通信保障、资源折中四份已验证预案分别执行完整分区审计，每份均有80个资源剖面交叉检查，全部一致。不同上游预案使用各自的真实日程与实际关联，不共用一套手工填入的时间区间。

\subsection{独立执行方案与资源缺口}
\subsubsection{两组独立执行的配置结果}
最优两组为$B_1\cup B_2$与$B_3$，即除S006外14区和S006独立。需求与缺口为
\begin{equation}
 \boldsymbol D^{(2)}=(4,2,3,6,4,4,2,3),\qquad
 \boldsymbol G^{(2)}=(0,0,1,0,0,0,0,0).
\label{eq:q4_two}
\end{equation}
总需求28，小于库存总件数30，但仍缺1架C型运输机。剩余中继组件不能替代飞机，故总件数并非可行性判据。工作量变异系数0.941762大于另两种分区，这是以较低资源缺口为优先产生的均衡代价，不应隐藏。

\subsubsection{三组独立执行的配置结果}
三组分区唯一，为$B_1|B_2|B_3$，得到
\begin{equation}
 \boldsymbol D^{(3)}=(5,2,5,7,4,6,2,3),\qquad
 \boldsymbol G^{(3)}=(1,0,3,1,0,2,0,0).
\label{eq:q4_three}
\end{equation}
需补1架A型运输机、3架C型运输机、1组A型电池和2组C型电池，共7件；总配置34，工作量变异系数1.183798。虽然$34-30=4$，仍有3组中继组件剩余，正缺口应为$4+3=7$，不能跨类型净额相抵。图\ref{fig:q4_resources}统一展示各类型库存与需求。
\fig[.9]{q4_resources.pdf}{固定工期优先预案后的分类型资源配置。格内为整数件数，红框表示超过同类型库存；不同资源不能按总件数相互替代。}{fig:q4_resources}

\subsubsection{分组引起的资源共享损失}
对每种资源都有
\begin{equation}
 \sum_g\max_t n_{gp}(t)\ge\max_t\sum_gn_{gp}(t).
\label{eq:q4_pooling}
\end{equation}
左侧为各独立组峰值之和，右侧为全局共享峰值。不同组峰值若错时，共享资源可以先后接续；独立后每组须为自己的峰值配置。C型运输机的共享、两组与三组需求分别为2、3、5架，任务时刻不变，增加来自禁止跨组复用。并非所有类型都增加：两组C型电池仍为4组，中继实体及组件分别保持2、3。

\subsection{稳健性与敏感性分析}
\label{sec:q4_robust}
\subsubsection{分类型库存增减与关键资源识别}
逐类型将部署前库存减少1或增加1，连同基准共17种库存，对两组与三组分别重新评价全部合法分区，形成34条汇总及完整候选明细。表\ref{tab:q4_inventory_stress}给出正缺口响应。增加一架C型运输机可使两组缺口从1降为0；增加一组中继能源组件不改变缺口。减少一组中继组件也不改变结果，因为该类型需求3、库存6，但这不是对任意设备执行中故障的保证。
\begin{table}[htbp]\centering\small
\caption{各类型库存单项增减后的最小总缺口}\label{tab:q4_inventory_stress}
\begin{tabular}{@{}lrrrr@{}}
\toprule
扰动资源 & 两组：减1 & 两组：加1 & 三组：减1 & 三组：加1\\
\midrule
A型运输机 & 2 & 1 & 8 & 6\\
B型运输机 & 2 & 1 & 8 & 7\\
C型运输机 & 2 & 0 & 8 & 6\\
A型电池 & 2 & 1 & 8 & 6\\
B型电池 & 2 & 1 & 8 & 7\\
C型电池 & 2 & 1 & 8 & 6\\
中继无人机 & 2 & 1 & 8 & 7\\
中继能源组件 & 1 & 1 & 7 & 7\\
\bottomrule
\end{tabular}
\par\smallskip\parbox{.96\linewidth}{\footnotesize 每次只改变一类库存，其余保持原值；两组三组重新比较全部合法分区。基准缺口分别为1与7。}
\end{table}


两组三组在该资源减少情景下的差异，来自原库存与峰值需求之间的结构余量。试验指部署前可用件数变化，允许重新指派同型资源，不模拟某个正在执行任务的设备突然失效。后者需要任务恢复和通信重算，超出当前固定日程试验。

\subsubsection{库存变化下的分组选择稳定性}
17种库存条件下，两组最优分区均不改变。完整枚举还揭示比有限网格更强的原因：另外两种分区的需求均为$(5,2,4,7,4,6,2,3)$，逐分量不小于原选定分区$(4,2,3,6,4,4,2,3)$，且总件数33大于28。

因此，对任意非负库存向量$\boldsymbol I'$都有
\begin{equation}
 \sum_p\max(D_p^{\mathrm{selected}}-I_p',0)
 \le\sum_p\max(D_p^{\mathrm{other}}-I_p',0).
\label{eq:q4_inventory_dominance}
\end{equation}
若第一目标相等，第二目标仍由28对33严格胜出。所以在日程、不可拆关系及既定词典序目标均固定时，单纯改变库存不会改变该两组选择。该结论是由完整候选向量得到的分量支配推论，不是把17个情景的经验稳定性推广到所有扰动；改变工作量优先关系或重新优化日程后不再有同一保证。

\subsubsection{不同运输—中继预案的配置代价}
进一步将四份问题三候选分别冻结，再独立重建关联和资源峰值。它们的工期、能源与分组需求见表\ref{tab:q4_crossplans}。工期优先方案的两组缺口1、三组7；其余三份两组均为2，节能方案三组仍为7，增强保障与资源折中方案三组为6。
\begin{table}[htbp]\centering\small
\caption{不同上游预案冻结后的独立配置比较}\label{tab:q4_crossplans}
\begin{tabular}{@{}lrrrrr@{}}
\toprule
预案 & 工期/s & 能耗/kWh & 附加损耗/dB & 两组缺口 & 三组缺口\\
\midrule
工期优先 & 5836.970 & 68.939416 & 0.00 & 1 & 7\\
节能备选 & 5911.548 & 68.735791 & 0.00 & 2 & 7\\
增强通信保障 & 5866.141 & 69.031751 & 0.25 & 2 & 6\\
三组资源折中 & 5866.141 & 68.975497 & 0.00 & 2 & 6\\
\bottomrule
\end{tabular}
\end{table}


其中资源折中方案比工期优先方案慢29.171\,s，三组少补1件，但两组多缺1件。由此不能以“资源更省”概括其所有组织目标；还需指明组数与类型。比较中每个点对应不同上游预案，不是保持同一日程时通过调节组数来改变执行工期。跨组复制中继或任意平移任务将改变题设边界，本章不混用这些反事实结果。
